% Source-derived chapter 01: Introduction
% Original source: hitchin-fibrations-abelian-surfaces-pw
\chapter{Introduction}
\label{hitchin-fibrations-abelian-surfaces-pw-chapter-01}
\source{hitchin-fibrations-abelian-surfaces-pw}

\begin{remark}[Source section]
\label{hitchin-fibrations-abelian-surfaces-pw-chapter-01-source}
\source{hitchin-fibrations-abelian-surfaces-pw}
\source{hitchin-fibrations-abelian-surfaces-pw}
This chapter reproduces the corresponding section of the retrieved
LaTeX source, including its definitions, statements, examples, and
proofs. It has not yet been translated into Lean.
\notready
\end{remark}

% The source section title was: Introduction
\subsection{Perverse filtrations}
 Throughout the paper, we work over the complex numbers $\BC$.
 
Let $\pi: X \rightarrow Y$ be a proper morphism with $X$ a nonsingular algebraic variety. The perverse $t$-structure on the constructible derived category $D_c^b(Y)$ with rational coefficients $\BQ$ induces a finite and  increasing filtration on the rational cohomology groups $H^*(X, \BQ)$
\begin{equation} \label{Perv_Filtration}
\source{hitchin-fibrations-abelian-surfaces-pw}
    P_0H^\ast(X, \BQ) \subset P_1H^\ast(X, \BQ) \subset \dots \subset P_kH^\ast(X, \BQ) \subset \dots \subset H^\ast(X, \BQ),
\end{equation}
called the \emph{perverse filtration} associated with $\pi$. The filtration (\ref{Perv_Filtration}) is governed by the topology of the morphism $\pi: X\rightarrow Y$. See Section \ref{sec1} for a brief review of the subject.

Every cohomology class on $X$ is indexed by the perverse filtration (\ref{Perv_Filtration}). We say that a class $\alpha \in H^d(X, \BQ)$ has \emph{perversity} $k$ if $\alpha = 0$, or
\[
\alpha \in P_kH^d(X, \BQ) \; \mathrm{and} \;  \alpha \notin P_{k-1}H^d(X, \BQ).
\]


The purpose of this paper is to study perverse filtrations associated with Hitchin fibrations in view of the \emph{P=W conjecture} by de~Cataldo, Hausel, and Migliorini \cite{dCHM1}. Our method is to use \emph{symmetries} induced by the monodromy of moduli spaces of sheaves on abelian surfaces.

\subsection{The P=W conjecture} 
Let $C$ be an irreducible  nonsingular  projective curve of genus $g \geq 2$. There are two moduli spaces which are attached to the curve $C$, the reductive Lie group $\mathrm{GL}_r$, and an integer $\chi$ with $\mathrm{gcd}(r,\chi)=1$.  
They are the twisted versions of Simpson's Dolbeault and Betti moduli spaces \cite{Si1994II}; see \cite{HT1} for  the non-abelian Hodge theory in the twisted  case.


The first moduli space $\CM_{\mathrm{Dol}}$ parametrizes stable Higgs bundles on $C$
\[
(\CE, \theta), \quad \theta: \CE \to \CE \otimes \Omega_C
\]
with $\mathrm{rank}(\CE)=r$ and $\chi(\CE) =\chi$. The variety $\CM_{\mathrm{Dol}}$ admits a projective morphism with connected fibers
\begin{equation}\label{Hitchin_fib}
\source{hitchin-fibrations-abelian-surfaces-pw}
h: \CM_{\mathrm{Dol}} \rightarrow \Lambda = \bigoplus_{i=1}^n H^0(C, \Omega_C^{\otimes i})
\end{equation}
sending $(\CE, \theta) \in \CM_{\mathrm{Dol}}$ to the characteristic polynomial $\mathrm{char}(\theta) \in \Lambda$. The proper morphism (\ref{Hitchin_fib}), called the \emph{Hitchin fibration}, is Lagrangian with respect to the canonical holomorphic symplectic form on~$\CM_{\mathrm{Dol}}$ given by the hyper-K\"ahler metric on $\CM_{\mathrm{Dol}}$; see \cite{Hit1}.  
The second moduli space is the  (twisted) character variety $\CM_B$. It  can be described (see \cite{HT1}, or  \cite{Shende} for a an alternative description) as parametrizing isomorphism classes of irreducible local systems
\[
\rho: \pi_1(C \setminus \{p\}) \to \mathrm{GL}_r
\]
where $\rho$ sends a loop around a chosen point $p$ to $\xi^\chi_r I_r \in \mathrm{GL}_r$, with $\xi_r$ a primitive root of unity. The character variety $\CM_B$ is affine. 

In \cite{Simp}, Simpson constructed a diffeomorphism between the (un-twisted) moduli spaces 
$\CM_{\mathrm{Dol}}$ and $\CM_B$, called the non-abelian Hodge theory;  see \cite{HT1} for the twisted case.  A striking 
prediction, suggested by the parallel between the Relative Hard Lefschetz \cite{BBD} and 
Curious Hard Lefschetz \cite{HRV} Theorems, was formulated by de~Cataldo, Hausel, and Migliorini \cite{dCHM1}; it predicts that the perverse filtration of $\CM_{\mathrm{Dol}}$ with respect to the Hitchin fibration (\ref{Hitchin_fib}) matches the weight filtration of the mixed Hodge structure on $\mathcal{M}_B$ under the identification
\[
H^\ast(\CM_{\mathrm{Dol}}, \BQ) = H^\ast(\CM_B, \BQ) 
\]
induced by Simpson's  (twisted) non-abelian Hodge theory.

\begin{conj}[\cite{dCHM1} P=W]
\source{hitchin-fibrations-abelian-surfaces-pw}
\notready \label{P=W_conj}
\source{hitchin-fibrations-abelian-surfaces-pw}
We have
\[
P_kH^\ast(\CM_{\mathrm{Dol}}, \BQ) = W_{2k}H^\ast(\CM_{B}, \BQ)= W_{2k+1}H^\ast(\CM_{B}, \BQ), \quad k \geq 0.
\]
\end{conj}

Conjecture \ref{P=W_conj} establishes a surprising connection between the topology of Hitchin fibrations and the Hodge theory of character varieties. It was proven in \cite{dCHM1} in the case of $r=2$ for any genus $g \geq 2$. See also \cite{dCHM3, SZ, Z} for certain parabolic cases, and \cite{SY, HLSY} for a compact analog concerning Lagrangian fibrations on projective holomorphic symplectic manifolds. 

The first main result of this paper is a proof of Conjecture \ref{P=W_conj} for curves of genus $2$ and arbitrary rank $r\geq 1$.

\begin{thm}
\source{hitchin-fibrations-abelian-surfaces-pw}
\notready\label{genus_2}
\source{hitchin-fibrations-abelian-surfaces-pw}
The P=W
Conjecture \ref{P=W_conj} holds when $C$ has genus $g=2$.
\end{thm}

For an irreducible nonsingular curve $C$ of genus $g \geq 2$, it is a general fact that $P=W$ for $\mathrm{GL}_r$ implies $P=W$ for $\mathrm{PGL}_r$; see (\ref{eqn71}) and (\ref{eqn89}). Hence we conclude immediately from Theorem \ref{genus_2} that the $P=W$ conjecture holds for $\mathrm{PGL}_r$ when the curve $C$ has genus $g=2$.

As explained in \cite[Section 1]{dCHM1}, the curious Poincar\'e duality and the Curious Hard Lefschetz conjectures \cite[Conjectures 4.2.4 and 4.2.7]{HRV} for character varieties are consequences of the P=W Conjecture \ref{P=W_conj}.  Moreover, by \cite{CDP} and \cite[Section 9.3]{MT}, the P=W Conjecture \ref{P=W_conj} implies the correspondence between Gopakumar--Vafa invariants and Pandharipande--Thomas invariants \cite[Conjecture 3.13]{MT} for the local Calabi--Yau 3-fold $T^*C \times \BC$ in the curve class $r[C]$. Hence Theorem \ref{genus_2} verifies all these conjectures for genus 2 curves.\footnote{As we discuss later, the curious Poincar\'e and Lefschetz conjectures have been recently shown by Mellit \cite{Mellit} to hold for all genera.}

%The right-hand side of (\ref{eqP=W}) has been studied in \cite{HLR, HRV}. Its closed formula was conjectured in \cite[Conjecture 1.2.1]{HLR} in terms of combinatorial data.

%Furthermore, the perverse numbers $^\mathfrak{p}h^{i,j}(\CM_{\mathrm{Dol}})$ appear naturally as refined Gopakumar--Vafa invariants, which count curves on the local Calabi--Yau~$3$-fold
%\[
%T^\ast C \times \BC;
%\]
%see \cite{CDP, MT} and Section \ref{sec0.4.1} for more discussions.

\subsection{Tautological classes}\label{sec_0.3}
\source{hitchin-fibrations-abelian-surfaces-pw}
Assume that $C$ has genus $g \geq 2$ and $r \geq 1$. A set of generators for the  cohomology rings identified by the non-abelian Hodge diffeomorphism
\begin{equation}\label{cohomology}
\source{hitchin-fibrations-abelian-surfaces-pw}
H^\ast(\CM_{\mathrm{Dol}}, \BQ) = H^\ast(\CM_B, \BQ) 
\end{equation}
is described in \cite{Markman} by tautological classes. 

Let
\[
p_C: C \times \CM_{\mathrm{Dol}} \to C,\quad  p_\CM: C \times \CM_{\mathrm{Dol}} \to  \CM_{\mathrm{Dol}}\]
be the projections. We say that a triple
\[
(\CU^\alpha, \theta) = (\CU, \theta, \alpha)
\]
is a {\em twisted universal family} over $C \times \CM_{\mathrm{Dol}}$, if $(\CU, \theta)$ is a universal family and 
\[
\alpha = p_C^\ast \alpha_1 +p_\CM^\ast \alpha_2 \in H^2(C\times \CM_{\mathrm{Dol}}, \BQ),
\]
with $\alpha_1 \in H^2(C, \BQ)$ and $\alpha_2 \in H^2(\CM_{\mathrm{Dol}}, \BQ)$.

For a twisted universal family $(\CU^\alpha, \theta)$, we define the {\em twisted Chern character} $\mathrm{ch}^\alpha(\CU)$ as
\begin{equation*}\label{uni_eqn1}
\source{hitchin-fibrations-abelian-surfaces-pw}
\mathrm{ch}^\alpha(\CU) = \mathrm{ch}(\CU) \cup \mathrm{exp}(\alpha) \in H^*(C\times \CM_{\mathrm{Dol}}, \BQ),
\end{equation*}
and we denote by 
\[
\mathrm{ch}_k^\alpha(\CU) \in H^{2k}(C\times \CM_{\mathrm{Dol}},\BQ)
\]
its degree $k$ part. The class $\mathrm{ch}^\alpha(\CU)$ is called \emph{normalized} if 
\[
\mathrm{ch}^\alpha_1(\CU)|_{p \times \CM_{\mathrm{Dol}}} =0 \in H^2(\CM_{\mathrm{Dol}}, \BQ), \qquad \mathrm{ch}^\alpha_1(\CU)|_{C \times q} =0 \in H^2(C, \BQ),
\]
with $p \in \CM_{\mathrm{Dol}}$ and $q\in C$ points.  Since two universal families differ by the pull-back
of a line bundle on $\CM_{\mathrm{Dol}}$,  a straightforward calculation shows that  normalized classes on $C\times \CM_{\mathrm{Dol}}$  exist
and are uniquely determined. We introduce the tautological classes associated with the normalized class $\mathrm{ch}^\alpha(\CU)$ as follows. 


%\footnote{A universal family $(\CU, \theta)$ on $C\times \CM_{\mathrm{Dol}}$ may not be deformed when we deform the curve $C$, but deformations of the Chern character $\mathrm{ch}(\CU)$ are always twisted Chern characters. Hence it is more natural to work with twisted universal families and twisted Chern characters for any curve $C$.}


For any $\gamma \in H^i(C, \BQ)$, let $c(\gamma ,k)$ denote the \emph{tautological class}
\begin{equation}\label{taut_class}
\source{hitchin-fibrations-abelian-surfaces-pw}
c(\gamma ,k) := \int_\gamma \mathrm{ch}_k^\alpha(\CU) = {p_{\CM}}_\ast( p_C^\ast \gamma \cup \mathrm{ch}^\alpha_k(\CU)) \in H^{i+2k-2}(\CM_{\mathrm{Dol}}, \BQ).
\end{equation}
Markman showed in \cite{Markman} that the classes $c(\gamma,k)$ generate the cohomology ring (\ref{cohomology}) as a $\BQ$-algebra. Furthermore, Shende proved in \cite{Shende} { (cf. Lemma \ref{lem3.1})}  that 
\begin{equation}\label{taut_weight}
\source{hitchin-fibrations-abelian-surfaces-pw}
    c(\gamma,k) \in {^k\mathrm{Hdg}^{i+2k-2}(\CM_B)},\quad \forall~\gamma \in H^i(C, \BQ),
\end{equation}
where 
\[
^k\mathrm{Hdg}^d(\CM_B) =  W_{2k}H^d(\CM_B, \BQ) \cap F^k H^d(\CM_B, \BC) \cap \bar{F}^k H^d(\CM_B, \BC),
\]
with $F^*$ and $W_*$ the Hodge  and weight filtrations, respectively. It follows that  the rational cohomology  $H^*(\CM_B, \BQ)$ is split  of Hodge--Tate type, and there is a canonical decomposition of 
the graded vector spaces (\ref{cohomology}) induced by the mixed  Hodge theory of $\CM_B$,
\begin{equation}\label{weight_decomp}
\source{hitchin-fibrations-abelian-surfaces-pw}
H^\ast(\CM_{\mathrm{Dol}}, \BQ) = H^\ast(\CM_B, \BQ) = \bigoplus_{k,d}  {^k\mathrm{Hdg}^d(\CM_B)}.
\end{equation}

%\footnote{By \cite{Shende}, every class in $H^2(\CM_B, \BQ)$ lies in $^1\mathrm{Hdg}^2(\CM_B)$, therefore (\ref{taut_weight}) holds only for any $\alpha \in H^2(\CM_B, \BQ)$.} 

The P=W conjecture, 
which can be restated as predicting that
\begin{equation}\label{p=w-explicit}
\source{hitchin-fibrations-abelian-surfaces-pw}
P_k H^\ast(\CM_{\mathrm{Dol}}, \BQ) = \bigoplus_{k' \leq k}  {^{k'}\mathrm{Hdg}^*(\CM_B)},
\end{equation} 
 is equivalent to the following statements.

%\[
%P_k H^d(\CM_{\mathrm{Dol}}, \BQ) = \bigoplus_{i \leq k} G_iH^d(\CM_{\mathrm{Dol}}, \BQ)
%\]

\begin{conj}[Equivalent version of P=W]
\source{hitchin-fibrations-abelian-surfaces-pw}
\notready\label{P=W2} 
\source{hitchin-fibrations-abelian-surfaces-pw}
There exists a splitting $G_*H^*(\CM_{\mathrm{Dol}}, \BQ)$ of the perverse filtration associated with the Hitchin fibration (\ref{Hitchin_fib}) satisfying the following two properties.
\begin{enumerate}
    \item[(a)](Tautological classes) 
   Every tautological class $c(\gamma,k)$ has perversity $k$ and in fact
    \begin{equation*}\label{taut_conj}
\source{hitchin-fibrations-abelian-surfaces-pw}
    c(\gamma,k) \in G_kH^{\ast}(\CM_{\mathrm{Dol}}, \BQ),\quad \forall k \geq 0, \, \forall \gamma \in H^*(C, \BQ).
    \end{equation*}
  
    \item[(b)] (Multiplicativity) The perverse decomposition is multiplicative, \emph{i.e.}, 
    \[
    \cup: G_k H^d(\CM_\mathrm{Dol}, \BQ)\times G_{k'} H^{d'}(\CM_\mathrm{Dol}, \BQ) \to G_{k+k'} H^{d+d'}(\CM_\mathrm{Dol}, \BQ).
    \]
\end{enumerate}
\end{conj}

We recall that the multiplicativity of the weight filtration
\[
\cup: W_k H^d(X, \BQ) \times W_{k'} H^{d'}(X, \BQ) \to W_{k+k'}H^{d+d'}(X, \BQ)
\]
is standard from mixed Hodge theory. However, the perverse filtration associated with a proper flat morphism is not multiplicative in general; see \cite[Exercise 5.6.8]{Park}. It is mysterious why the perverse filtration associated with the Hitchin fibration (\ref{Hitchin_fib}) should be multiplicative, as is predicted by the P=W conjecture. In fact, one consequence of this paper is that the full P=W conjecture can be reduced to the multiplicativity of the perverse filtration; see Theorem \ref{last}.  This approach of analyzing the P=W conjecture via Conjecture \ref{P=W2} goes back to \cite{dCHM1} where it is applied to prove the case of rank $2$.

\begin{comment}f
The Curious Hard Lefschetz \cite[Conjecture 4.2.7]{HRV}, which corresponds to the relative hard Lefschetz for the Hitchin fibration (\ref{Hitchin_fib}), is expected to be induced by the class 
\begin{equation}\label{rel_ample}
\source{hitchin-fibrations-abelian-surfaces-pw}
c_\CU([\mathrm{pt}], 0) \in  H^2(\CM_B, \BQ) = H^2(\CM_\mathrm{Dol}, \BQ).
\end{equation}
Since it is relatively ample with respect to the Hitchin morphism, the splitting $G_*H^*(\CM_{\mathrm{Dol}}, \BQ)$ of the perverse filtraion in Conjecture \ref{P=W2} should also be induced by the class (\ref{rel_ample}).

For any proper morphism, Deligne constructed in \cite{D} three splittings of the perverse filtration via a fixed relative ample class. Then de Cataldo \cite{dC} extended Deligne's constructions and gave two additional splittings. All these five splittings may be different in general.

Now if we choose a universal family $(\CU, \theta)$, we obtain the tautological classes (\ref{taut_class}) as well as five (possibly distinct) splittings
\begin{equation}\label{perv_split}
\source{hitchin-fibrations-abelian-surfaces-pw}
H^*(\CM_{\mathrm{Dol}}, \BQ) = \bigoplus_{i,d}G_iH^d(\CM_{\mathrm{Dol}}, \BQ)
\end{equation}
induced by the class (\ref{rel_ample}).

The following theorem proves Conjecture \ref{P=W2} (a) for all these five splittings. Moreover, it shows that the multiplicativity holds for the subring of $H^\ast(\CM_{\mathrm{Dol}}, \BQ)$ generated by even tautological classes.
\end{comment}

%We note that the classes (\ref{taut_class}) are dependent on the choice of the universal bundle $\CU$, and the difference of two universal bundles is given by the pullback of a line bundle $\CL$ on $\CM_{\mathrm{Dol}}$,
%\[
%\CU'= \CU \otimes {p_\CM}^\ast \CL.
%\]
%However, since (\ref{taut_weight}) holds for any choice of universal family, it suffices to prove Conjecture \ref{P=W2} for only \emph{one} universal family. 

%The following theorem proves Conjecture \ref{P=W2} (a). It also verify the multiplicativity for the sub-ring of $H^\ast(\CM_{\mathrm{Dol}}, \BQ)$ generated by even tautological classes.

\begin{comment}
The existence of the decomposition
\[
H^*(\CM_{\mathrm{Dol}}, \BQ) = \bigoplus_{k,d} G_k H^d(\CM_{\mathrm{Dol}}, \BQ) 
\]
proposed in Conjecture \ref{P=W2} is predicted by the decomposition (\ref{weight_decomp}), which is uniquely characterized by the properties (a) and (b) in Conjecture \ref{P=W2}. 
\end{comment}

\subsection{P=W for tautological classes}
Theorem \ref{genus_2} is concerned with genus $g=2$.
In this section, we state our results for genus $g \geq  2$.
\subsubsection{Even tautological classes}
We consider the $\BQ$-subalgebra
\[
R^*(\CM_{\mathrm{Dol}})  \subset H^*(\CM_{\mathrm{Dol}}, \BQ)= H^*(\CM_B, \BQ)
\]
generated by all the \emph{even} tautological classes
\[
c(\gamma,k), \quad  k\geq 0, \quad  \gamma \in H^0(C, \BQ)\oplus H^2(C, \BQ).
\]
Our next theorem establishes the P=W conjecture in arbitrary genus and rank, but restricted to this subalgebra.
\begin{thm}
\source{hitchin-fibrations-abelian-surfaces-pw}
\notready\label{P=W_even}
\source{hitchin-fibrations-abelian-surfaces-pw}
The P=W conjecture holds, for any genus $g \geq 2$ and rank $r\geq 1$ for $R^*(\CM_{\mathrm{Dol}})$, \emph{i.e.}
\[
\begin{split}
P_kH^\ast(\CM_{\mathrm{Dol}}, \BQ) \cap R^*(\CM_{\mathrm{Dol}}) & = W_{2k}H^\ast(\CM_{B}, \BQ) \cap R^*(\CM_{\mathrm{Dol}})\\
& = W_{2k+1}H^\ast(\CM_{B}, \BQ) \cap R^*(\CM_{\mathrm{Dol}}).
\end{split}
\]
\end{thm}

In fact, in the proof of Theorem  \ref{P=W_even}, we obtain a multiplicative decomposition 
\begin{equation*}
R^*(\CM_{\mathrm{Dol}}) = \bigoplus_{k,d} G_k R^d(\CM_{\mathrm{Dol}})
\end{equation*}
 that splits the restricted perverse filtration $P_kH^\ast(\CM_{\mathrm{Dol}}, \BQ) \cap R^*(\CM_{\mathrm{Dol}})$ and such that 
\[
G_k R^d(\CM_{\mathrm{Dol}}) =  {^k\mathrm{Hdg}^d}(\CM_B) \cap R^{d}(\CM_{\mathrm{Dol}}).
\]
We refer to Section \ref{Section4.6} for more details.

\begin{comment}
Theorem \ref{P=W_generator} (a) reduces Conjecture \ref{P=W_conj} to the multiplicativity (Conjecture \ref{P=W2} (b)) for arbitrary one of the five splittings induced by (\ref{rel_ample}). We expect that all these five splittings coincide. This was proven for $n=2$ in \cite{dCHM1}. The equivalence of the five splittings also hold for a Lagrangian fibration associated with a projective holomorphic symplectic variety; see \cite[Remark 1.3]{SY}.
\end{comment}

\subsubsection{Odd tautological classes}
The following theorem concerns odd tautological classes
\begin{equation}\label{odd_taut}
\source{hitchin-fibrations-abelian-surfaces-pw}
c(\gamma, k) \in H^{2k-1}(\CM_{\mathrm{Dol}}, \BQ), \quad \forall k \geq 0, \, \forall \gamma \in H^1(C, \BQ).
\end{equation}

\begin{thm}
\source{hitchin-fibrations-abelian-surfaces-pw}
\notready\label{P=W_odd}
\source{hitchin-fibrations-abelian-surfaces-pw}
Any odd tautological class (\ref{odd_taut}) has perversity $k$.
\end{thm}

Theorems \ref{P=W_even} and \ref{P=W_odd} provide strong evidence for the P=W Conjecture
\ref{P=W_conj} (and for its reformulation Conjecture \ref{P=W2}) for any genus $g \geq 2$. In particular, we obtain that all tautological generators $c(\gamma,k)$ satisfy the $P=W$ match:
\begin{equation}\label{eqn9}
\source{hitchin-fibrations-abelian-surfaces-pw}
    c(\gamma, k) \in P_k H^*(\CM_{\mathrm{Dol}}, \BQ),\quad \forall~ \gamma \in H^*(C, \BQ).
\end{equation}

\subsubsection{Multiplicativity is equivalent to P=W} The P=W Conjecture \ref{P=W_conj}  implies immediately that the perverse filtration $P_*H^*(\CM_{\mathrm{Dol}}, \BQ)$ is multiplicative, \emph{i.e.}
\begin{equation}\label{multi1}
\source{hitchin-fibrations-abelian-surfaces-pw}
\cup: P_k H^d(\CM_\mathrm{Dol}, \BQ)\times P_{k'} H^{d'}(\CM_\mathrm{Dol}, \BQ) \to P_{k+k'} H^{d+d'}(\CM_\mathrm{Dol}, \BQ).
\end{equation}
The following theorem shows that the converse is also true. It is a corollary of Theorems \ref{P=W_even} and \ref{P=W_odd}, and of the Curious Hard Lefschetz conjecture recently established by Mellit \cite{Mellit}.

\begin{thm}
\source{hitchin-fibrations-abelian-surfaces-pw}
\notready\label{last}
\source{hitchin-fibrations-abelian-surfaces-pw}
The P=W Conjecture \ref{P=W_conj} is equivalent to the multiplicativity (\ref{multi1}) of the perverse filtration.
\end{thm}

In fact, we deduce from (\ref{taut_weight}), (\ref{eqn9}), and (\ref{multi1}) that
\[
W_{2i}H^*(\CM_B, \BQ) \subset P_iH^*(\CM_{\mathrm{Dol}}, \BQ), \quad \forall i\geq 0,
\]
which further implies Conjecture \ref{P=W_conj} by \cite[Theorem 1.5.3]{Mellit}
and the second paragraph following \cite[Corollary 1.5.4]{Mellit}.


{
\begin{rmk}
\source{hitchin-fibrations-abelian-surfaces-pw}
\notready\label{r=1}
\source{hitchin-fibrations-abelian-surfaces-pw}
All the results of this paper hold for arbitrary rank $r \geq 1$. The case $r=1$ is classical and easy: the Dolbeault moduli space is the cotangent bundle of a Jacobian of some degree  of the curve and the Betti moduli space
is a product of $\mathbb G_m^{2g}$. In this paper, starting with the assumption {$\beta^2\geq 6$} in Section \ref{Sec2.1}, we focus on the case of rank $r\geq 2$, {which corresponds to $\beta^2 \geq 8$}. The case of rank one could be dealt with also by using the methods of this paper, but at the price of some easy modifications that we deemed a distraction. 
\end{rmk}
}
\subsection{Idea of the proof}

The key idea in proving our results
is to first study perverse filtrations for certain \emph{compact} hyperk\"ahler manifolds, motivated in part by earlier work of the third author and Yin \cite{SY}.
That is, we consider an abelian surface $A$ with ample curve class $\beta \in H_2(A, \BZ)$, and prove analogs of 
our main result for the moduli space
$\CM_{\beta,A}$ of one-dimensional sheaves with support in the class $\beta$; see Theorem \ref{taut_abelian} for the precise statement.
The advantage of the compact geometry is that we can use
results of Markman \cite{Markman3} on monodromy operators for $\CM_{\beta,A}$.  These are symmetries, arising from parallel transport and Fourier-Mukai transforms, that do not appear in the Hitchin setting and which heavily constrain tautological classes of universal families. We show that these operators relate perverse filtrations and tautological classes for different choices of $\beta$ and, in particular, allow us to pass from the case of imprimitive $\beta$ to primitive $\beta$, which can be studied directly using \cite{Z,SZ}.

In order to apply this to Hitchin fibrations, we consider the degeneration to the normal cone of an embedding
of {a genus $g$ curve $C$} into an abelian surface
\[
j_C: C\hookrightarrow A,
\]
and study the specialization map on cohomology of the associated moduli spaces.  In general, there is a great deal of information loss in this specialization because the family is non-proper. Nevertheless, we are able to show its compatibility with tautological classes and perverse splittings.  Finally we use these compatibilities  to deduce our main theorems {for such a curve $C$; this implies our results hold  for all curves in view of \cite{dCM}.}

\subsection{Outline of paper}

We briefly outline the contents of this paper.  In Section $1$, we recall some basic facts about perverse filtrations and earlier results of \cite{Z,SZ} for Hilbert schemes of points on the special abelian surface $E \times E'$.
In Section $2$, we pass to the setting of abelian surfaces.  One  key result  here is the splitting Theorem \ref{taut_abelian}.  After recalling Markman's work on monodromy operators, we prove the result for primitive classes $\beta$ by a direct analysis, and then we combine the primitive case with Markman's results  to establish the general case.
In Section $3$, we formulate and prove  Theorem \ref{strengthened2.1},  which is a strengthened version of the splitting of Theorem \ref{taut_abelian} and which  is  more robust for specialization arguments.
Finally, in section $4$, we study  specialization maps for our degeneration, and we use them to prove the main theorems.

While this paper is written in logical order, some readers may prefer to start from \S\ref{se4}, and refer backwards as needed.




\subsection{Acknowledgements}
We are grateful to Tam\'as Hausel, Jochen Heinloth, Luca Migliorini, Vivek Shende, David Vogan, Qizheng Yin, Zhiwei Yun, and Zili Zhang for helpful discussions.  M.A. de Cataldo is partially supported by NSF grants DMS 1600515 and 1901975.  D. Maulik is partially supported by NSF FRG grant DMS-1159265. J. Shen is supported by the NSF grant DMS 2134315.
